\section{Background}
\label{sec:background}

\subsection{The layout problem}
\label{sec:layout-problem}
A layout is the order of a tensor's dimensions in memory, for example \{1,0,2\} rather than \{0,1,2\}.
In TPUGraphs, XLA chooses layouts only at configurable nodes (convolution, dot and reshape), whose TPU
implementations are the most constrained, and derives every other node's layout from those
choices~\cite{phothilimthana2023tpugraphs}. A layout configuration is one set of choices across all
configurable nodes of a graph. The choices interact: when two connected operations use different
layouts, the compiler inserts a copy (a transpose) between them, so nodes cannot be optimised
independently and a whole-graph model is needed. XLA's autotuner finds 10--20\% speedups on production
models, but every candidate must be compiled and timed. A learned cost model therefore ranks the
candidates from the graph so that only the top few are timed~\cite{kaufman2021learned}. The graph is
the HLO input to XLA's layout-assignment step: nodes are tensor operations and edges are data flow.

\subsection{The four collections}
\label{sec:collections}
TPUGraphs' layout data come in four collections: two graph sources crossed with two sampling schemes.
The xla graphs come from the XLA regression benchmark and MLPerf, covering vision, NLP, speech and
recommendation; the nlp graphs are BERT variants that differ in layers, attention heads and hidden
size. Random sampling draws diverse candidates with widely spread runtimes, while default sampling runs
a genetic search from the compiler's default layout and yields similar candidates. Default is harder
than random~\cite{phothilimthana2023tpugraphs}, and we find xla harder than nlp
(Section~\ref{sec:result1}). Google measured every runtime on one TPU v3 core. The xla collections have
7 validation graphs each and the nlp collections 20; the released test files carry no runtimes, so
every number in this paper is on validation graphs.

\subsection{Related work}
\label{sec:related}
Learned cost models span tensor-program models (AutoTVM~\cite{chen2018learning},
TenSet~\cite{zheng2021tenset}, TLP~\cite{zhai2022deep}), basic-block throughput models
(Ithemal~\cite{mendis2019ithemal}, GRANITE~\cite{sykora2022granite}, with the analytical
Facile~\cite{abel2023facile} as counterpoint), whole-graph models on TPU and XLA (Kaufman et
al.~\cite{kaufman2021learned}, TPUGraphs~\cite{phothilimthana2023tpugraphs},
TGraph~\cite{khizbullin2024graph}, Polynormer~\cite{deng2024polynormer},
PerfSAGE~\cite{chai2023perfsage}) and program representations (ProGraML~\cite{cummins2021programl}).
Recent directions include NeuSight~\cite{lee2024forecasting}, compiler world
models~\cite{pan2026toward}, regression language models~\cite{akhauri2025regression} and
Omniwise~\cite{wang2025omniwise}. These works report progress as one model beating another. We do not
propose a new model; we test how much any static-graph model can know, and Section~\ref{sec:audit}
audits how such comparisons are tested.
